\documentclass[../../../../main.tex]{subfiles}

\begin{document}

\section*{第1問}

\begin{proof}[解]
\[
a(n)=\frac{(n+2)(n+3)(n+4)}{n!}
\]

\textbf{(1)} $a(n)>0$であり，
\[
\frac{a(n+1)}{a(n)}=\frac{(n+3)(n+4)(n+5)/(n+1)!}{(n+2)(n+3)(n+4)/n!}=\frac{n+5}{(n+1)(n+2)}\longrightarrow0 \quad(n\to\infty)
\]
となり，$n$が十分大きい所で比が$1/2$未満になることから，$a(n)\to0$である。

\textbf{(2)} $a(1)$から書き出す。
\[
a(1)=a(2)=60
\]
\[
a(3)=35,\quad a(4)=14,\quad a(5)=\frac{21}{5},\quad a(6)=1
\]
\[
a(7)<1
\]
以下，$a(n)$が$n\ge2$で単調減少であることを示す。
\[
a(n)-a(n+1)=\frac{(n+3)(n+4)}{(n+1)!}\left[(n+1)(n+2)-(n+5)\right]
\]
\[
=\frac{(n+3)(n+4)}{(n+1)!}(n^2+2n-3)>0 \quad(\because n\ge2)
\]
より，示された。これと$a(7)<1$，$a_n\to0\ (n\to\infty)$から，$7\le n$なる$n\in\mathbb{N}$に対し，$0<a(n)<1$となり，$a(n)\notin\mathbb{Z}$。以上から，
\[
n=1,2,3,4,6
\]

\textbf{(3)} (2)から，$S_n=\displaystyle\prod_{k=1}^na(k)$とおくと，$n\ge7$の時，$S_n$は単調減少である。(2)から，$n=1,2,\cdots,7$の時，$S_n\in\mathbb{Z}$である。
\[
S_7=a(1)\cdots a(7)=2^2\cdot3^3\cdot5^2\cdot7^2\cdot11=1455300\in\mathbb{Z}
\]
$n\ge8$の時
\[
S(8)=S_7\cdot a(8)=\frac{3^2\cdot5^2\cdot7\cdot11^2}{2^2}\notin\mathbb{Z}
\]
\[
S(9)=S(8)\cdot a(9)=\frac{5\cdot11^3\cdot13}{2^7\cdot3}\notin\mathbb{Z}
\]
\[
S(10)=S(9)\cdot a(10)=\frac{11^3\cdot13^2}{2^{12}\cdot3^4\cdot5}\notin\mathbb{Z}<1
\]
より，$10\le n$の時，$0<S(n)<1$で$S(n)\notin\mathbb{Z}$。以上から，
\[
n=1,2,3,4,5,6,7
\]
\end{proof}

\newpage

\section*{第2問}

\begin{proof}[解]
5点$A,B,C,D,E$を，$A$を原点，$AB$を$x$軸正方向にとり，$AB=BC=CD=DE=1$，各頂点での外角を$\theta$として配置する。すなわち
\[
A(0,0),\ B(1,0),\ C(1+\cos\theta,\sin\theta),
\]
\[
D(1+\cos\theta+\cos2\theta,\ \sin\theta+\sin2\theta),
\]
\[
E(1+\cos\theta+\cos2\theta+\cos3\theta,\ \sin\theta+\sin2\theta+\sin3\theta)
\]
とおける。五角形$ABCDE$の面積$S$は，座標による面積公式（シューレースの公式）から
\[
2S=\sum(x_iy_{i+1}-x_{i+1}y_i)
\]
であり，$A\to B$，$E\to A$の寄与は0，$B\to C$の寄与は$\sin\theta$，$C\to D$の寄与は
\[
(1+\cos\theta)(\sin\theta+\sin2\theta)-(1+\cos\theta+\cos2\theta)\sin\theta=\sin2\theta+(\cos\theta\sin2\theta-\cos2\theta\sin\theta)=\sin2\theta+\sin\theta
\]
$D\to E$の寄与は
\[
(1+\cos\theta+\cos2\theta)\sin3\theta-(\sin\theta+\sin2\theta)\cos3\theta=\sin3\theta+\sin2\theta+\sin\theta
\]
だから，
\[
2S=3\sin\theta+2\sin2\theta+\sin3\theta
\]
\[
S=\frac12\sin3\theta+\sin2\theta+\frac32\sin\theta
\]
\[
\frac{dS}{d\theta}=\frac32\cos3\theta+2\cos2\theta+\frac32\cos\theta
\]
\[
=\frac32(\cos3\theta+\cos\theta)+2\cos2\theta=3\cos2\theta\cos\theta+2\cos2\theta=\cos2\theta(3\cos\theta+2)
\]
$0<\theta<\pi/2$から$\cos\theta\in(0,1)$で$3\cos\theta+2>0$だから，$dS/d\theta$の符号は$\cos2\theta$と一致し，下表を得る。
\begin{center}
\begin{tabular}{c|ccccc}
$\theta$ & $0$ & & $\pi/4$ & & $\pi/2$ \\
\hline
$S'$ & & $+$ & $0$ & $-$ & \\
\hline
$S$ & & $\nearrow$ & & $\searrow$ &
\end{tabular}
\end{center}
従って，$S$は$\theta=\pi/4$の時，最大値
\[
S\left(\frac{\pi}{4}\right)=\frac32\cdot\frac{\sqrt2}{2}+1+\frac12\cdot\frac{\sqrt2}{2}=\sqrt2+1
\]
をとる。
\end{proof}

\newpage

\section*{第3問}

\begin{proof}[解]
\textbf{(1)} $f(x)=\dfrac{x^2}{n^2}+e^{2x}-1$，$f'(x)=\dfrac{2}{n^2}x+2e^{2x}$。$f'(x)$は単調増加から，$f'(x)\to\pm\infty\ (x\to\pm\infty)$から，$f'(\alpha)=0$なる$\alpha$がただ1つ存在。（$\alpha<0$）
\begin{center}
\begin{tabular}{c|ccccc}
$x$ & & & $\alpha$ & & \\
\hline
$f'$ & & $-$ & $0$ & $+$ & \\
\hline
$f$ & & $\searrow$ & & $\nearrow$ &
\end{tabular}
\end{center}
（$f(x)\to\infty\ (x\to\pm\infty)$，$f(0)=0$で，グラフは下図の概形。）

\textbf{(2)} $x_n$は$\dfrac{1}{n^2}x_n^2+e^{2x_n}-1=0$をみたし，$x_n\ne0$であるから，$f(x)=0$の$x\ne0$の解である。
\[
f(-n)=e^{-2n}>0 \quad\cdots\text{①}
\]
$k$を定数とする（$k>0$）。
\[
f(-n+k)=\frac{-2kn+k^2}{n^2}+e^{2(k-n)}
\]
\[
=-2k\cdot\frac1n+\frac{k^2}{n^2}+e^{2(k-n)}
\]
$n$が十分大きい時，$\dfrac1n\gg\dfrac1{n^2}\gg\dfrac1{e^n}$だから
\[
f(-n+k)<0 \quad\cdots\text{②}
\]
①，②と$y=f(x)$のグラフから
\[
-n<x_n<-n+k
\]
$n>0$から
\[
-1<\frac{x_n}{n}<-1+\frac{k}{n}
\]
はさみうちの定理から
\[
\frac{x_n}{n}\longrightarrow-1 \quad(n\to\infty)
\]
\end{proof}

\newpage

\section*{第4問}

\begin{proof}[解]
$n\ge3\ \cdots$①

\textbf{(1)} 全ての場合$nP_2$通りについての引いた2枚の積の和は（1枚目，2枚目で区別）
\[
1\cdot(2+\cdots+n)+2(1+3+\cdots+n)+\cdots+n(1+2+\cdots+n-1)
\]
\[
=(1+2+\cdots+n)^2-(1^2+2^2+\cdots+n^2)
\]
\[
=\frac14n^2(n+1)^2-\frac16n(n+1)(2n+1)
\]
だから，
\[
E=\frac{1}{{}_nP_2}\left[\frac14n^2(n+1)^2-\frac16n(n+1)(2n+1)\right]
\]
\[
=\frac{1}{n-1}\left[\frac14n(n+1)^2-\frac16(n+1)(2n+1)\right]
\]
\[
=\frac{(n+1)(3n^2-n-2)}{12(n-1)}=\frac{(n+1)(3n+2)}{12}
\]

\textbf{(2)} 全ての場合$nP_3$通りについての引いた3枚の積は
\[
(1+2+\cdots+n)^3-{}_3C_2(1^2+2^2+\cdots+n^2)(1+2+\cdots+n)+2(1^3+2^3+\cdots+n^3)
\]
\[
=\left\{\frac12n(n+1)\right\}^3-3\cdot\frac12n(n+1)\cdot\frac16n(n+1)(2n+1)+2\left\{\frac12n(n+1)\right\}^2
\]
\[
=\left\{\frac12n(n+1)\right\}^2\left[\frac12n(n+1)-(2n+1)+2\right]
\]
だから，
\[
\frac{E(n)}{n^3}=\frac{\frac14n^2(n+1)^2\left(\frac12n^2-\frac32n+1\right)}{n^4(n-1)(n-2)}
\]
\[
=\frac14\left(1+\frac1n\right)^2\cdot\frac{(n-1)(n-2)}{2(n-1)(n-2)}\longrightarrow\frac18 \quad(n\to\infty)
\]
\end{proof}

\end{document}
